\documentclass[11pt]{article}
\usepackage[margin=1in]{geometry}
\usepackage[T1]{fontenc}
\usepackage[utf8]{inputenc}
\usepackage{lmodern,microtype,amsmath,amssymb,amsthm,mathtools}
\usepackage[colorlinks=true,linkcolor=blue,citecolor=blue,urlcolor=blue]{hyperref}
\hypersetup{pdftitle={Exponential online mistakes for isotropic Bayesian triangle-wave networks},pdfauthor={ChatGPT-6 Astra; Kaarel Hanni; Dmitry Vaintrob}}
\newtheorem{theorem}{Theorem}
\newtheorem{lemma}[theorem]{Lemma}
\newtheorem{proposition}[theorem]{Proposition}
\newtheorem{corollary}[theorem]{Corollary}
\newtheorem{remark}[theorem]{Remark}
\newcommand{\R}{\mathbb R}
\newcommand{\F}{\mathbb F}
\newcommand{\E}{\mathbb E}
\newcommand{\Pp}{\mathbb P}
\newcommand{\one}{\mathbf 1}
\newcommand{\sgn}{\operatorname{sign}}
\newcommand{\cV}{\mathcal V}
\newcommand{\eps}{\varepsilon}
\title{Exponential online mistakes for isotropic Bayesian\protect\\triangle-wave networks}
\author{ChatGPT-6 Astra \and Kaarel H\"anni \and Dmitry Vaintrob\thanks{Authors are listed in alphabetical order.}}
\date{7 November 2026}
\begin{document}
\maketitle

\begin{abstract}
We give a fixed continuous, bounded, $1$-Lipschitz activation for which
an isotropic Gaussian prior on network weights makes
superpolynomially many online mistakes on a fixed input permutation.
The activation is the period-four triangle wave. For every $m\ge5$,
the target has $d=\Theta(m^6)$ Boolean inputs, a bounded-fan-in Boolean
circuit of size $\Theta(d)$, and an explicit one-hidden-layer teacher
of width $2m-1$ with every weight at most one and output margin one.
At every sufficiently large learner width, taking the Gaussian variance
to zero at that fixed width makes the actual posterior probability of
the correct label tend to zero at each of $2^{m-4}$ designated rounds.
The learner conditions on absolute output margin greater than $1/2$
on the full Boolean cube before conditioning on exact labels.
The proof identifies an exact weighted Fourier variation norm,
excludes competing sign functions outside an explicitly defined parity
family by a quantitative margin argument, and strictly orders the
remaining alternatives, each wrong on one logical input class, using
distinct square roots. Integer neuron allocations transfer the order to every
sufficiently large finite width, and Gaussian energy separation controls
the posterior mixture. The mistake bound is exponential in the displayed
teacher width and is $\exp(\Omega(d^{1/6}))$ in input dimension and
Boolean-circuit size. Calibration occupies $15/16$ of the cube, and the
sufficient width is not claimed to be polynomial. The theorem does not
establish the corresponding statement for tanh, ReLU, GELU, or hardtanh.
\end{abstract}

\section{Model, target, and theorem}

Let $\phi$ be the odd triangle wave of period four, defined by
\[
 \phi(t)=t\quad(-1\le t\le1),\qquad
 \phi(t+2)=-\phi(t).
\]
Thus $\phi$ is continuous, $1$-Lipschitz, bounded by one, and vanishes
at zero. The same activation is used at every value of the construction
parameter. On $X=\{-1,1\}^d$, write $\bar x=(1,x)$ and use one hidden
layer,
\begin{equation}\label{eq:model}
 F_{W,v}(x)=\sum_{a=1}^N v_a\phi(W_a\cdot\bar x),\qquad
 W_{ai}\sim N(0,N\sigma^2/d),\quad
 v_a\sim N(0,\sigma^2/N),
\end{equation}
independently, including the bias coordinate $i=0$. There is no
coordinate-dependent covariance. Equivalently, the whitened coordinates
$\xi=\sqrt{d/N}\,W$ and $\eta=\sqrt N\,v$ are independent centered
Gaussians with common variance $\sigma^2$. The Gaussian energy is
\begin{equation}\label{eq:energy}
 K_N(W,v)=\frac dN\|W\|_F^2+N\|v\|_2^2.
\end{equation}
Before observing labels the learner conditions on
\begin{equation}\label{eq:globalmargin}
 \mathcal M=\{|F_{W,v}(x)|>1/2\text{ for every }x\in X\}.
\end{equation}
It subsequently conditions on exact noiseless training signs. All
posterior statements below refer to this actual conditional measure,
not to a selected minimum-energy network.
The bounded-weight teacher complexity $C_2(f)$ is the least integer
$C\ge1$ for which a width-$C$ network in the same unnormalized predictor
form as \eqref{eq:model}, with all incoming and readout weights at most
$C$ in absolute value, implements $f$ with margin greater than $1/2$.

Fix $m\ge5$, and set
\begin{equation}\label{eq:dimensions}
 q_i=2m+2i-1,\quad r_i=2q_i\quad(1\le i\le m),\quad
 D=\sum_i r_i=6m^2,\quad
 c=10^6(D+1)^3+1,\quad d=c+D.
\end{equation}
Partition the coordinates into a common core $C$ of odd size $c$ and
blocks $B_i$ of even size $r_i$. For a set $A$, write
$\chi_A(x)=\prod_{j\in A}x_j$. Define $z(x)\in G:=\F_2^m$ by
\[
 (-1)^{z_i(x)}=\chi_{B_i}(x).
\]
Our target is
\begin{equation}\label{eq:target}
 f_m(x)=\chi_C(x)\bigl(2\one_{\{z(x)=0\}}-1\bigr).
\end{equation}

\begin{theorem}\label{thm:main}
For every $m\ge5$, the target \eqref{eq:target} and a fixed permutation
of the full cube have the following properties.
\begin{enumerate}
\item The target has a bounded-fan-in Boolean circuit of size
$\Theta(d)$ and an explicit one-hidden-layer triangle-wave teacher of
width $2m-1$, with all incoming and readout weights bounded by one and
absolute output margin at least one.
In particular, $C_2(f_m)\le2m-1$.
\item The permutation has a calibration prefix of length
$15\cdot2^{d-4}$, followed by $2^{m-4}$ designated, nonrepeated inputs.
There is $N_0(m)$ such that for every fixed $N\ge N_0(m)$ and every
designated round,
\begin{equation}\label{eq:posteriorlimit}
 \lim_{\sigma\downarrow0}
 \Pp_{N,\sigma}\bigl(F(x)f_m(x)>0\mid\mathcal M,
       \text{all preceding exact labels}\bigr)=0.
\end{equation}
The permutation is independent of $N$ and $\sigma$.
\item At each such fixed width, posterior-label majority makes all
$2^{m-4}$ designated mistakes for all sufficiently small $\sigma$;
the posterior wrong-label probability also exceeds $2/3$ there.
The probability that posterior sampling makes every designated mistake
tends to one.
\end{enumerate}
The number of designated mistakes is exponential in the displayed
teacher width, and is $\exp(\Omega(d^{1/6}))$. The theorem asserts no
variance threshold uniform in width and no polynomial bound on $N_0(m)$.
\end{theorem}

\begin{corollary}\label{cor:width}
For the explicit unit-bounded teacher widths $w_m=2m-1$, the construction
forces $2^{(w_m-7)/2}$ designated mistakes, at every sufficiently large
learner width in the fixed-width strong-prior limit. Its input dimension
and bounded-fan-in Boolean-circuit size are $\Theta(w_m^6)$.
\end{corollary}

The calibration and width qualifications are substantive. The result
rules out an arbitrary-order polynomial mistake bound in the target's
Boolean-circuit size, but does not give $2^{\Omega(d)}$ mistakes or a
short calibration prefix. All proofs are conventional mathematical
proofs; no formal verification is claimed.

\section{An exact weighted Fourier norm}

For a cube function $F$, let $\widehat F(A)=\E F(X)\chi_A(X)$.
For $s\in G$, set
\[
 A_s=C\cup\bigcup_{i:s_i=1}B_i,\qquad
 a_s=\sqrt{|A_s|}=\sqrt{c+\sum_i r_i s_i}.
\]
The code projection is the orthogonal projection onto these characters:
\begin{equation}\label{eq:projection}
 PF(x)=\sum_{s\in G}\widehat F(A_s)\chi_{A_s}(x)
      =\chi_C(x)\,\bar F(z(x)).
\end{equation}
Equivalently,
$\bar F(z)=\E[\chi_C(X)F(X)\mid z(X)=z]$.
Define
\[
 Q(F)=\sum_{s\in G}a_s|\widehat F(A_s)|,
 \qquad
 \|F\|_{\cV}=\inf\left\{\sum_a|v_a|\|\theta_a\|_2:
          F=\sum_a v_a\phi(\theta_a\cdot\bar x)\right\},
\]
where representations are finite and equality is on the cube. This
variation cost is finite on every cube function: each odd-support
Fourier character is a pure atom with bias zero, each even-support
character is a pure atom with bias one, and the constant uses
$\phi(1)=1$. Homogeneity and the triangle inequality follow by scaling
and concatenating representations, while
$\|F\|_\infty\le\sqrt{d+1}\|F\|_{\cV}$ proves positivity. Thus it is
a norm.

\begin{theorem}[Atomic norm and quantitative rounding]\label{thm:atomic}
More generally, let $c\ge1000$ be odd, let the disjoint blocks have even
sizes at least four and total size $D$, and assume
$c\ge100(D+1)^2$. For every biased triangle-wave ridge
$h_\theta(x)=\phi(\theta\cdot\bar x)$,
\begin{equation}\label{eq:atomic}
 Q(h_\theta)\le\|\theta\|_2.
\end{equation}
Nonzero equality occurs exactly at the minimal pure-character atoms:
zero bias, weights $\pm1$ on the core and on a union of whole blocks,
and zero weights elsewhere. Consequently,
\begin{equation}\label{eq:exactnorm}
 \left\|\sum_{s\in G} b_s\chi_{A_s}\right\|_{\cV}
       =\sum_{s\in G}a_s|b_s|.
\end{equation}
Moreover, with $\kappa=1/200$, every atom has a number $e_\theta\ge0$
such that
\begin{equation}\label{eq:rounding}
 \|\theta\|_2-Q(h_\theta)\ge\kappa e_\theta,
\end{equation}
and the following two bounds hold:
\begin{align}
 |h_\theta(x)-h_\theta(x')|
   &\le e_\theta\|x-x'\|_2
   &&\text{if }(\chi_C,z)(x)=(\chi_C,z)(x'),\label{eq:fiberlip}\\
 |\E_{X_C}h_\theta(X_C,u)|
   &\le\sqrt{D+1}\,e_\theta
   &&\text{for every fixed physical block assignment }u.
   \label{eq:coremean}
\end{align}
\end{theorem}

The exactness and the quantitative statement are both needed. An
approximate norm identity need not preserve the small strict energy
differences used later. A proof, including constants, is given in
Appendix~\ref{app:atomic}.

For odd $|A|$, the pure atom is especially simple:
\begin{equation}\label{eq:pureparity}
 \phi\!\left(\sum_{i\in A}x_i\right)
       =(-1)^{(|A|-1)/2}\chi_A(x).
\end{equation}
Thus all witnesses in \eqref{eq:exactnorm} use ordinary bias-free
triangle-wave neurons with incoming weights zero or one, with a sign
absorbed into the readout.

\section{Global margin excludes non-code sign completions}

The projection $P$ alone does not preserve global absolute margin.
The next lemma is the step which prevents fractional projected
completions from invalidating the argument.

\begin{lemma}[A margin-aware gap]\label{lem:offcode}
Under the assumptions of Theorem~\ref{thm:atomic}, suppose
$|F(x)|\ge1/2$ everywhere. If its sign is not of the form
\begin{equation}\label{eq:codesign}
 \sgn F(x)=\chi_C(x)g(z(x)),\qquad g:G\to\{-1,1\},
\end{equation}
then
\begin{equation}\label{eq:offcodegap}
 \|F\|_{\cV}\ge Q(F)+\frac{\kappa}{2\sqrt{D+1}}.
\end{equation}
\end{lemma}
\begin{proof}
For any representation $F=\sum_a v_a h_{\theta_a}$, put
$V=\sum_a|v_a|\|\theta_a\|_2$ and
$B=\sum_a|v_a|e_{\theta_a}$. The Fourier triangle inequality and
\eqref{eq:rounding} give
\begin{equation}\label{eq:VQgap}
 V-Q(F)\ge\sum_a|v_a|\bigl(\|\theta_a\|_2-Q(h_{\theta_a})\bigr)
       \ge\kappa B.
\end{equation}
Each fiber of $(\chi_C,z)$ is connected by moves which flip two
coordinates within one block, with the core regarded as a block.
If the sign is nonconstant on a fiber, some such edge has opposite
endpoint signs. The margin gives an output difference at least one,
whereas \eqref{eq:fiberlip} bounds it by $2\sqrt2 B$. Hence
$B\ge1/(2\sqrt2)$.

Otherwise the sign is constant on every fiber. Failure of
\eqref{eq:codesign} means that at some logical control value both core
parity classes have the same sign. Fix physical block coordinates with
that control value. Every core assignment then has that sign and
margin, so $|\E_{X_C}F|\ge1/2$. Equation~\eqref{eq:coremean} gives
$B\ge1/(2\sqrt{D+1})$. Since $D\ge4$, this latter lower bound suffices
in either case. Take the infimum over representations in
\eqref{eq:VQgap}.
\end{proof}

Let $M=2^m$, and define the affine free set
\begin{equation}\label{eq:freecoset}
 H=\{z:z_1=z_2=z_3=z_4=0\},\qquad
 U=e_1+H,\qquad L=H\cup U.
\end{equation}
In particular $|U|=M/16$ and $0\notin U$. The calibration reveals
the target on every physical input with $z(x)\notin U$. It has
exactly $15\cdot2^{d-4}$ inputs.

\begin{lemma}[The flat calibration bound]\label{lem:flat}
Every real $f:G\to\R$ with
$f(0)\ge1/2$ and $f(z)\le-1/2$ for
$z\notin\{0\}\cup U$ satisfies
\begin{equation}\label{eq:flatlower}
 \sum_s|\widehat f(s)|\ge L_0:=\frac32-\frac4M.
\end{equation}
For every $v\in U$, the star score
\begin{equation}\label{eq:starscore}
 f_v=-\frac12+\one_{\{0\}}+\one_{\{v\}}
\end{equation}
attains equality and has global margin $1/2$.
\end{lemma}
\begin{proof}
Consider
\[
 \mathcal D(f)=\left(1-\frac4M\right)f(0)
  -\frac4M\sum_{z\in H\setminus\{0\}}f(z)
  -\frac2M\sum_{z\notin L}f(z).
\]
Its Fourier multipliers are $-1$ at zero, $+1$ on
$L^\perp\setminus\{0\}$, $3/4$ on
$H^\perp\setminus L^\perp$, and $+1$ outside $H^\perp$.
Thus $\mathcal D(f)\le\sum_s|\widehat f(s)|$. The training inequalities
give $\mathcal D(f)\ge3/2-4/M$. Finally,
\[
 \widehat f_v(0)=-\frac12+\frac2M,
 \qquad
 \widehat f_v(s)=\frac{1+(-1)^{s\cdot v}}M\quad(s\ne0),
\]
whose absolute sum is $L_0$.
\end{proof}

\begin{proposition}[Reduction after calibration]\label{prop:codereduction}
After the calibration, and after any subsequent labels for which at
least one star remains consistent, every minimum-variation sign
completion is of the form \eqref{eq:codesign}. Every non-code sign
completion has variation cost strictly larger than every star cost.
\end{proposition}
\begin{proof}
For every training-consistent $F$, its projection \eqref{eq:projection}
satisfies the latent calibration margins, because these were imposed on
complete physical fibers. Lemma~\ref{lem:flat} gives
\[
 Q(F)\ge\sqrt c\,L_0.
\]
A non-code sign completion therefore costs at least
\begin{equation}\label{eq:noncodebarrier}
 B_{\rm nc}:=\sqrt c\,L_0+\frac{\kappa}{2\sqrt{D+1}}.
\end{equation}
By \eqref{eq:exactnorm}, every star has cost at most
\begin{equation}\label{eq:starupper}
 \sqrt{c+D}\,L_0
   \le\sqrt c\,L_0+\frac{3D}{4\sqrt c}
   <B_{\rm nc},
\end{equation}
where the strict inequality follows from \eqref{eq:dimensions} and
$\kappa=1/200$. Later labels retain the initial lower bound and a
remaining star upper witness. If a sign completion is odd and
code-invariant, its projection $PF$ has the same signs and global
margin, and \eqref{eq:exactnorm} gives a code score of no larger cost.
Thus the remaining optimization is the finite weighted Fourier margin
problem.
\end{proof}

\section{The strictly ordered latent completions}

Write $w(s)=a_s$ and
\[
 B_w(f)=\sum_s w(s)|\widehat f(s)|,
 \qquad
 \widehat w(v)=\frac1M\sum_s w(s)(-1)^{s\cdot v}.
\]
The weights satisfy $w(s)\ge w(0)=\sqrt c$.

\begin{lemma}[Signs and distinctness of the weight transform]\label{lem:weights}
Every nonconstant Walsh coefficient $\widehat w(v)$ is strictly
negative, and all $M$ values $\widehat w(v)$ are distinct.
\end{lemma}
\begin{proof}
The integral representation
\[
 \sqrt x=\frac1{2\sqrt\pi}\int_0^\infty(1-e^{-xt})t^{-3/2}\,dt
\]
gives, for $v\ne0$,
\[
 \widehat w(v)=-\frac1{2\sqrt\pi M}
 \int_0^\infty e^{-ct}
 \prod_{i:v_i=1}(1-e^{-r_i t})
 \prod_{i:v_i=0}(1+e^{-r_i t})\,t^{-3/2}\,dt<0.
\]
The integral is convergent at zero because $v\ne0$.

For distinctness, all distinct integers in $[c,c+D]$ have different
squarefree parts. Indeed, if $n=qu^2<n'=qv^2$ have the same squarefree
part, then
$n'-n\ge q(2u+1)>2\sqrt c>D$, a contradiction. Square roots of
distinct squarefree integers are rationally linearly independent.
Consequently the distinct values $\sqrt n$, $c\le n\le c+D$, are
rationally linearly independent.

The coefficients of $M\widehat w(v)$ as a linear combination of these
square roots are the coefficients of
\[
 P_v(t)=\prod_{i:v_i=1}(1-t^{r_i})
         \prod_{i:v_i=0}(1+t^{r_i}).
\]
The odd integers $q_i$ lie strictly between $2m$ and $4m$, so none
divides another. The cyclotomic polynomial $\Phi_{q_i}$ divides
$1-t^{2q_j}$ exactly when $i=j$, and never divides
$1+t^{2q_j}$: a root of odd order cannot have a power equal to $-1$.
Thus $\Phi_{q_i}$ divides $P_v$ exactly when $v_i=1$. All the
polynomials $P_v$ are distinct, proving the claim.
\end{proof}

\begin{lemma}[Subgroup classification costs]\label{lem:subgroup}
If a proper subgroup $J\le G$ has codimension at least two, the unique
minimum-$B_w$ real score with margin $1/2$ and positive set exactly $J$
is $f_J=\one_J-1/2$. In particular, writing $C_0$ for the singleton
cost and $C_v$ for the cost of $\{0,v\}$,
\begin{align}
 C_0&=\frac12\sqrt c+\widehat w(0)-\frac{2\sqrt c}M,\label{eq:singletoncost}\\
 C_v&=\frac12\sqrt c+\widehat w(0)+\widehat w(v)-\frac{4\sqrt c}M.
 \label{eq:starcost}
\end{align}
All $C_v$ are distinct and each is strictly smaller than $C_0$.
\end{lemma}
\begin{proof}
A norm-dual function has Fourier multipliers $+w(s)$ on
$J^\perp\setminus\{0\}$, $-w(0)$ at zero, and zero elsewhere. In
physical coordinates it is
\[
 \lambda_J(x)=\sum_{s\in J^\perp}w(s)(-1)^{s\cdot x}-2w(0).
\]
On $J$ it is at least $(|J^\perp|-2)w(0)>0$. For $x\notin J$,
Fourier inversion gives
\[
 \sum_{s\in J^\perp}w(s)(-1)^{s\cdot x}
   =\frac M{|J|}\sum_{h\in J}\widehat w(x+h)<0,
\]
so $\lambda_J(x)<0$. Its normalized pairing with a feasible score is
therefore minimized exactly when all signed margins equal $1/2$.
That score also saturates the Fourier dual inequality, proving
optimality and uniqueness. Direct evaluation gives
\eqref{eq:singletoncost}--\eqref{eq:starcost}. The final claims follow
from Lemma~\ref{lem:weights}, since
$C_0-C_v=2\sqrt c/M-\widehat w(v)>0$.
\end{proof}

\begin{lemma}[Non-subgroup gap]\label{lem:nonsubgroup}
Suppose a globally margin-$1/2$ score $f$ has a positive set $A$
containing zero, and let $\delta=|A|/M<1/3$. If $A$ is not a subgroup,
then
\[
 \sum_s|\widehat f(s)|\ge\frac{2-\delta}{1+2\delta}.
\]
For every calibration-consistent positive set this lower bound is at
least $61/38$.
\end{lemma}
\begin{proof}
Put $B=\sum_s|\widehat f(s)|$ and $a=\widehat f(0)$. Since
$|f(x)|\le B$,
$a\le\delta B-(1-\delta)/2$. Choose $u,v\in A$ with $u+v\notin A$.
Every nonconstant character takes the value $+2$ or $-2$ in
$\chi(0)+\chi(u)+\chi(v)-\chi(u+v)$, while the constant character
gives $+2$. Hence
\[
 2\le f(0)+f(u)+f(v)-f(u+v)
   \le2a+2(B-|a|).
\]
If $a<0$, combining these inequalities gives
$(1+2\delta)B\ge2-\delta$. If $a\ge0$, the mean bound gives
$B\ge(1-\delta)/(2\delta)$, which is at least the claimed bound
because $\delta<1/3$. Calibration allows positives only in
$\{0\}\cup U$, so $\delta\le1/16+1/M\le3/32$, giving $61/38$.
\end{proof}

Thus every non-subgroup code completion has variation cost at least
\begin{equation}\label{eq:nonsubgroupbarrier}
 B_{\rm ns}:=\frac{61}{38}\sqrt c
    >\frac32\sqrt{c+D},
\end{equation}
which is larger than every star cost. A subgroup contained in
$\{0\}\cup U$ has at most two elements: two distinct elements of $U$
would sum to a nonzero element of $H$, which is forbidden. The only
possible subgroup positive sets are therefore the singleton and the
stars. These statements remain true when some points of $U$ have
subsequently been labeled negative.

\section{A fixed sequence and strict variation gaps}

Order the points of $U$ uniquely by
\[
 C_{v_1}<C_{v_2}<\cdots<C_{v_K},\qquad K=|U|=2^{m-4}.
\]
First present every calibration input in any fixed order. Next, for
$j=1,\ldots,K$, present one fixed physical representative $x_j$ with
$z(x_j)=v_j$ and $\chi_C(x_j)=+1$. These inputs were absent from the
calibration, are distinct, and all have true label $-1$. Complete the
sequence to a permutation of the remaining cube points in any fixed
order.

Before round $j$, the star at $v_j$ is consistent with every preceding
label and predicts $+1$ at $x_j$. Any odd code completion consistent
with the history has positive set contained in
$\{0\}\cup\{v_j,\ldots,v_K\}$. By the preceding results, among all
globally margin-valid completions the cheapest sign function is uniquely
the star at $v_j$. More quantitatively, there is $\Delta_j>0$ such that
every history-consistent score predicting the correct label at $x_j$
satisfies
\begin{equation}\label{eq:variationgap}
 \|F\|_{\cV}\ge C_{v_j}+\Delta_j.
\end{equation}
Explicitly, one may take
\[
 \Delta_j=\min\bigl(\{B_{\rm nc},B_{\rm ns},C_0\}
                  \cup\{C_{v_\ell}:j<\ell\le K\}\bigr)-C_{v_j}>0.
\]
No quantitative lower bound on these algebraic differences is needed.

Only the initial calibration uses complete fibers. Later individual
queries eliminate the corresponding code stars. The initial flat
lower bound continues to hold, and every remaining star remains a
global upper witness, so Proposition~\ref{prop:codereduction} still
excludes all non-code signs at every designated round.

\section{Every sufficiently large finite width and the posterior}

\begin{lemma}[Canonical energy and integer allocation]\label{lem:allocation}
For every width and every network,
\begin{equation}\label{eq:energylower}
 K_N(W,v)\ge2\sqrt d\sum_a |v_a|\|W_a\|_2
            \ge2\sqrt d\,\|F_{W,v}\|_{\cV}.
\end{equation}
For every fixed star score $\chi_C f_v$, there are networks at all
sufficiently large integer widths whose outputs equal that score and
whose energies converge to $2\sqrt d\,C_v$.
\end{lemma}
\begin{proof}
The lower bound is the arithmetic--geometric mean inequality applied
to each neuron's two energy terms. Write the exact pure-atom expansion
of the star as $\sum_s b_s\phi(u_s\cdot\bar x)$, absorbing phases into
$b_s$, with $\|u_s\|_2=a_s$. Define
\[
 p_s=\frac{|b_s|}{\sqrt d\,a_s}\quad(b_s\ne0),\qquad
 P_*:=\sum_s p_s\le\frac{3/2}{\sqrt{dc}}<1.
\]
At width $N$, take $n_s=\lceil Np_s\rceil$ replicas of atom $u_s$,
each with readout $b_s/n_s$, and set all other neurons to zero.
For all sufficiently large $N$, $\sum_s n_s\le N$. The score is
unchanged, and its energy is
\[
 \sum_s\left[d\frac{n_s}{N}a_s^2
                 +\frac{b_s^2}{n_s/N}\right]
 \longrightarrow2\sqrt d\sum_s|b_s|a_s=2\sqrt d\,C_v.
\]
This convergence is along all sufficiently large integer widths, not
only a subsequence.
\end{proof}

There are finitely many designated rounds. Equations
\eqref{eq:variationgap} and \eqref{eq:energylower}, together with
Lemma~\ref{lem:allocation}, therefore give $N_0(m)$ such that for
every $N\ge N_0(m)$ and every round $j$ there is a wrong-label star
witness with energy strictly smaller than the energy infimum of all
correct-label completions. Slightly multiplying its readout weights
by a factor greater than one creates strict global margins and
preserves this strict energy gap.
Enlarge $N_0(m)$ to at least $2m-1$; the padded target teacher then
keeps every posterior along the full permutation well-defined.

\begin{lemma}[Finite-dimensional Gaussian separation]\label{lem:laplace}
Fix a width $N$ and a history. Suppose the correct-label event within
the strict global-margin posterior has energy at least $b$, while a
wrong-label, history-consistent strict-margin parameter vector has
energy strictly below $b$. Then the posterior correct-label
probability tends to zero as $\sigma\downarrow0$.
\end{lemma}
\begin{proof}
Continuity and finiteness of the input cube give an open parameter
ball around the wrong witness contained in the denominator event.
Shrink it so its energy is at most $a<b$. In whitened coordinates
the Gaussian density is proportional to
$\sigma^{-p}\exp(-K_N/(2\sigma^2))$, with fixed finite
$p=N(d+2)$. The ball probability is at least a positive constant times
$\sigma^{-p}e^{-a/(2\sigma^2)}$.

Choose $\eta\in(0,1)$ with $(1-\eta)b>a$. A Gaussian exponential
moment bounds the raw probability of $K_N\ge b$ by
$\eta^{-p/2}e^{-(1-\eta)b/(2\sigma^2)}$. The ratio tends to zero.
The normalizing constant for the pre-data conditioning on
$\mathcal M$ cancels in this ratio. The argument controls the entire
posterior event and makes no choice among tied parameter minima.
\end{proof}

Apply Lemma~\ref{lem:laplace} at each designated history to obtain
\eqref{eq:posteriorlimit}. A union bound over the finitely many rounds
gives the posterior-sampling claim, without any independence
requirement on the prediction draws. Taking the minimum of the
finitely many small-variance thresholds gives the majority and
$2/3$ conclusions.

\section{Teacher complexity and interpretation}

The target has the explicit score
\begin{equation}\label{eq:teacher}
 T_m(x)=\chi_C(x)\left[\sum_{i=1}^m\chi_{B_i}(x)-(m-1)\right].
\end{equation}
The bracket is $+1$ when every block parity is positive, and is at
most $-1$ otherwise. Thus its sign is \eqref{eq:target} and its
absolute output is at least one. Expand \eqref{eq:teacher} as the
$m$ characters $\chi_{C\cup B_i}$ and $m-1$ copies of
$-\chi_C$. Each support has odd size. By \eqref{eq:pureparity}, this
is a width-$2m-1$ teacher with incoming weights zero or one, zero
biases, and readouts $\pm1$.

A Boolean circuit computes each block parity and the core parity,
tests whether all block parities are positive, and combines with the
core parity. With fan-in two this uses $O(d)$ gates. Every input is
relevant, so any fan-in-two circuit has at least $d-1$ gates; the
Boolean-circuit complexity is $\Theta(d)$. The displayed teacher has
$O(md)=O(m^7)$ scalar parameters, all bounded by one. Hence every
usual bounded same-depth teacher-size convention is polynomial in
$m$, while the mistake count is $2^{m-4}$.

Because $D=6m^2$ and $c=10^6(D+1)^3+1$, one has $d=\Theta(m^6)$.
The lower bound is therefore stretched exponential in input dimension
and Boolean-circuit size, and exponential in the explicit teacher
width. The construction uses a less customary periodic activation;
continuity alone is insufficient for a polynomial arbitrary-order
mistake theorem, but the corresponding questions for the usual
monotone activations remain distinct. The proof uses the strong-prior
limit at each fixed eligible width. It makes no assertion about a
fixed positive variance, a width-uniform small-variance threshold,
or the canonical width-first mean-field limit.

\appendix
\section{Proof of the atomic theorem}\label{app:atomic}

We give all quantitative ingredients of Theorem~\ref{thm:atomic}.
The absolutely convergent Fourier series of the triangle wave is
\begin{equation}\label{eq:triangleseries}
 \phi(t)=\sum_{n\ge1\ {\rm odd}}\alpha_n
      (-1)^{(n-1)/2}\sin(n\pi t/2),
 \qquad \alpha_n=\frac8{\pi^2n^2},\quad\sum_n\alpha_n=1.
\end{equation}
Adding an even integer to any incoming weight, including the bias,
changes the ridge only by a global sign on the cube. Reduce each
parameter modulo two into $[-1,1]$, decreasing its norm and preserving
all absolute Fourier coefficients. We first work with this reduced
vector, of norm $r$. If $r=0$, the ridge is zero and all assertions
are immediate with $e_\theta=0$; henceforth suppose $r>0$.

\subsection{Projection bounds and low norms}

For a harmonic and an odd-size support $A$, its absolute Fourier
coefficient is
\[
 |\cos(n\pi b/2)|
 \prod_{i\in A}|\sin(n\pi w_i/2)|
 \prod_{i\notin A}|\cos(n\pi w_i/2)|.
\]
For block $B_j$, put
$p_j=\prod_{i\in B_j}|\sin(n\pi w_i/2)|$ and
$q_j=\prod_{i\in B_j}|\cos(n\pi w_i/2)|$.
H\"older gives $p_j+q_j\le1$ for block size at least two.
Consequently, with
$\mathcal L=\sum_s|\widehat h_\theta(A_s)|$, one has
\begin{equation}\label{eq:Lfirstharmonic}
 \mathcal L\le1,\qquad
 \mathcal L\le1-\alpha_1(1-P),
\end{equation}
where $P$ is the first-harmonic product
\[
 P=|\cos(\pi b/2)|\prod_{i\in C}|\sin(\pi w_i/2)|
                         \prod_j(p_j+q_j).
\]

For $0\le u\le1$, the triangle series gives
\[
 \sum_n\alpha_n\sin^2(n\pi u/2)=u,\qquad
 \sum_n\alpha_n\sin^2(n\pi u/2)\cos^2(n\pi u/2)
                  =\tfrac12\min(u,1-u).
\]
It follows that
\[
 \sum_n\alpha_n\sin^4(n\pi u/2)
    =u-\tfrac12\min(u,1-u)\le\tfrac12(u+u^2).
\]
H\"older on the $c$ core factors and $c\ge4$ therefore imply,
with $t=r/\sqrt c$,
\[
 \mathcal L\le\frac1c\sum_{i\in C}\sum_n\alpha_n
              |\sin(n\pi w_i/2)|^c\le\frac{t+t^2}{2}.
\]
Thus for $r\le.9\sqrt c$,
\begin{equation}\label{eq:lownormgap}
 Q(h_\theta)\le.95\sqrt{1.01}\,r<.955r,
\end{equation}
since $D/c\le.01$.

\subsection{Localization near a pure atom}

Suppose $r>.9\sqrt c$ and $P>.86$. For a first-harmonic block write
$a=\max(p_j,q_j)$ and $b'=\min(p_j,q_j)$. Because its size is at
least four, generalized H\"older gives
$\sqrt a+\sqrt{b'}\le1$. Since $a+b'>.86$, it follows that
$a>.85$ and
\[
 b'\le(1-\sqrt a)^2\le.273(1-a)^2,
 \qquad -\log a\le1.05[-\log(a+b')].
\]
For the last elementary inequality, put $u=1-a\le.15$, use
$\log(1+.273u^2/a)\le.273u^2/a$ and
$-\log a\ge u+u^2/2$.

Choose the nearest integer vector $u$: bias even, core coordinates
odd, and every coordinate of a block odd when its sine product
dominates or even when its cosine product dominates. Reduced
parameters and the strict localization imply that this is a minimal
pure atom: its coordinates are zero or $\pm1$ as specified in the
theorem. Put $\delta=\theta-u$ and $\rho=\|\delta\|_2$. The dominant
products are products of $\cos(\pi\delta_i/2)$. Since
$\cos(\pi s/2)\le e^{-\pi^2s^2/8}$ for $|s|\le1$,
\[
 \rho^2\le\frac8{\pi^2}\,1.05(-\log P)<.1296,
 \qquad \rho<.36.
\]
If $\rho_j$ is the perturbation norm on one block, its nondominant
product is at most $(\pi^4/256)\rho_j^4$. Every dominant product is
at least $.85$. Therefore
\[
 \log P\le-\frac{\pi^2}{8}\rho^2
          +\frac{\pi^4}{256\cdot.85}\rho^4
          \le-1.17\rho^2.
\]
For $\rho\le.36$ this implies $1-P\ge\rho^2$. Since
$\alpha_1>.8$, \eqref{eq:Lfirstharmonic} yields
\begin{equation}\label{eq:localquadratic}
 \mathcal L\le1-.8\rho^2.
\end{equation}

Let the pure ridge be $\eta\chi_{A_T}$, let $a=\sqrt{|A_T|}$, and
put $a_{\max}=\sqrt{c+D}$ and
$\Delta=a_{\max}-a\le D/(2\sqrt c)\le.05$.
The $1$-Lipschitz property gives
$|\widehat h_\theta(A_T)|\ge1-\E|\delta\cdot\bar X|\ge1-\rho$.
Hence
\[
 Q(h_\theta)\le a\mathcal L+
       \Delta(\mathcal L-|\widehat h_\theta(A_T)|)
       \le a-.8a\rho^2+\Delta\rho.
\]
Using $r\ge a-\rho$, we obtain
\begin{equation}\label{eq:quadraticdeficit}
 r-Q(h_\theta)\ge .8a\rho^2-(1+\Delta)\rho.
\end{equation}

\subsection{A strict projection gap for an absolute linear form}

For an affine linear form $Z=\delta_0+\sum_i\delta_iX_i$ and
$\rho=\|\delta\|_2$, we claim
\begin{equation}\label{eq:absolutelemma}
 \sum_{\emptyset\ne S\subseteq[m]}
       \left|\widehat{|Z|}\!\left(\bigcup_{j\in S}B_j\right)\right|
       \le\E|Z|-\gamma\rho,
 \qquad
 \gamma=\frac2\pi\left(1-e^{-1/2}-\frac1{16}\right)>.21.
\end{equation}
Use
$|z|=(2/\pi)\int_0^\infty(1-\cos(tz))t^{-2}\,dt$.
At frequency $t$, put
$p_j=\prod_{i\in B_j}|\sin(t\delta_i)|$ and
$q_j=\prod_{i\in B_j}|\cos(t\delta_i)|$.
The absolute sum of nonempty block-character coefficients of the
cosine is bounded by
\[
 N_t=r_t\left[\prod_j(p_j+q_j)-\prod_jq_j\right],
\]
where $r_t$ is the absolute product of the bias cosine and cosine
factors outside the blocks. Since $p_j+q_j\le1$,
\[
 N_t\le1-|\cos(t\delta_0)\prod_i\cos(t\delta_i)|
 \le D_t:=1-\cos(t\delta_0)\prod_i\cos(t\delta_i).
\]
For $t\rho\le1$, all cosine factors are positive. The inequality
$\cos s\le e^{-s^2/2}$ gives
$D_t\ge(1-e^{-1/2})t^2\rho^2$. By AM--GM and block sizes at least
four,
\[
 N_t\le\sum_jp_j\le t^4\rho^4/16\le t^2\rho^2/16.
\]
Integrating the nonnegative difference $D_t-N_t$ over
$0\le t\le1/\rho$ proves \eqref{eq:absolutelemma}.

\subsection{Completing the norm and quantitative estimates}

If $\rho\le1.5/\sqrt c$, then
$\max_x|\delta\cdot\bar x|\le1.5\sqrt{(c+D+1)/c}<1.51<2$.
At an odd lattice peak the exact triangle identity is therefore
\[
 h_\theta(x)=\eta\chi_{A_T}(x)(1-|\delta\cdot\bar x|).
\]
Multiplication by $\chi_{A_T}$ permutes the block characters, and
\eqref{eq:absolutelemma} gives
\[
 Q(h_\theta)\le a[1-\E|Z|]+a_{\max}[\E|Z|-\gamma\rho]
    \le a-(\gamma a_{\max}-\Delta)\rho.
\]
Consequently,
\begin{equation}\label{eq:smallrho}
 r-Q(h_\theta)\ge(\gamma a_{\max}-\Delta-1)\rho\ge.15\rho.
\end{equation}
If instead $\rho\ge1.5/\sqrt c$, then
\eqref{eq:quadraticdeficit} gives the same lower bound $.15\rho$.

If $r>.9\sqrt c$ but $P\le.86$, then
\[
 \frac{r-Q(h_\theta)}r
 \ge1-\frac{[1-.14(8/\pi^2)]\sqrt{1.01}}{.9}>.010.
\]
Together with \eqref{eq:lownormgap}, this proves the atomic inequality
and gives strictness except at pure atoms. The pure atoms attain
equality directly, establishing \eqref{eq:exactnorm} and all equality
cases. Undoing the modulo-two reduction can only increase the norm.

For quantitative rounding, set $e_\theta=r$ in the two far cases
and $e_\theta=\rho$ in the localized case. The preceding estimates
imply \eqref{eq:rounding} with $\kappa=1/200$. In a far case the atom
is $r$-Lipschitz, and at fixed physical controls its affine argument
has square mean at most $(D+1)r^2$ over the core. This proves
\eqref{eq:fiberlip}--\eqref{eq:coremean} there.

In a localized case, periodicity gives the global identity
\[
 \phi((u+\delta)\cdot\bar x)
    =\eta\chi_{A_T}(x)\psi(\delta\cdot\bar x),
 \qquad \psi(t)=\phi(1+t).
\]
The function $\psi$ is even and $1$-Lipschitz. On a fixed
$(\chi_C,z)$ fiber the character is constant, proving
\eqref{eq:fiberlip}. Pairing either value of any one core coordinate
in the core average gives
\[
 |\E_{X_C}\phi((u+\delta)\cdot\bar X)|
       \le\min_{i\in C}|\delta_i|
       \le\rho/\sqrt c\le\sqrt{D+1}\,e_\theta,
\]
which proves \eqref{eq:coremean} and completes the theorem.

\end{document}
